\section{Ie\c{s}an--Saint Venant Formulation}

\subsection{Relaxed elastostatics}

Consider a prismatic, homogeneous, and isotropic linearly elastic body 
\(
\Body=\Section\times(0,L)
\),
where \(\Section\subset\mathbb{R}^2\) is a bounded (possibly multiply connected) cross-section
spanned by \(\{\FrameBasis_y,\FrameBasis_z\}\) and the axis is \(\FrameBasis\).
Set \(\bm{r}=y\,\FrameBasis_y+z\,\FrameBasis_z\) and
\(
\mathbf{1}=\FrameBasis\otimes\FrameBasis+\mathbf{1}_\Section
\)
with
\(
\mathbf{1}_\Section=\FrameBasis_y\otimes\FrameBasis_y+\FrameBasis_z\otimes\FrameBasis_z
\).
The lateral boundary is \(\BodyLateralSurface=\SectionBoundary\times(0,L)\).

\begin{remark}
The developments are frame-free (only right-handedness is assumed). For implementation it is convenient
either to take the first coordinate along \(\FrameBasis\) (so \(\FrameBasis\equiv\bm\partial_x\)) or the last;
both conventions are compatible with the formulas below.
\end{remark}

The classical Saint-Venant problems are not posed as classical elastostatic boundary value problems, but rather what is known as \emph{relaxed} elastostatic problems. This distinction (elaborate...)

The elastostatic field \((\bm u,\StressTensor)\) satisfies
\begin{equation}
\left\{
\begin{aligned}
\Div\,\StressTensor&=\mathbf{0} &\qquad\text{in }\Body\\
\StressTensor\,\bm n&=\mathbf{0} & \text{on }\BodyLateralSurface\\
\StressTensor&=\lambda\,(\operatorname{tr}\StrainTensor)\,\mathbf{1}+2G\,\StrainTensor,\\
\StrainTensor&=\tfrac12\big(\nabla\bm u+\nabla\bm u^{\!\top}\big),
\end{aligned}\right.,
\label{eq:partial-elastostatic-equilibrium}
\end{equation}
together with on \(\BodyLateralSurface\).
End tractions on \(\Section_0,\Section_L\) are \emph{not} prescribed pointwise; instead we constrain only the resultants
\begin{equation}
\mathscr{F}(\bm u)=\int_{\Section_1}\StressTensor(\bm u)\,(-\FrameBasis)\,\mathrm{d}A,\qquad
\mathscr{M}(\bm u)=\int_{\Section_1}\bm{r}\times\big(\StressTensor(\bm u)\,(-\FrameBasis)\big)\,\mathrm{d}A,
\label{eq:def-resultants}
\end{equation}
leaving the traction distributions on the end faces unknown.


At this stage we \emph{do not} impose any hypothesis on the in–plane stress block.  
We will track the axial normal stress \(\AxialStress:=\FrameBasis\cdot\StressTensor\,\FrameBasis\in\mathbb{R}\) and, when convenient,
a pointwise scalar shear component \(\PointShearStress\in\mathbb{R}\) together with the in–section shear vector \(\bm\tau\).
For bookkeeping it is useful to write the full decomposition
\[
\StressTensor=\AxialStress\,\FrameBasis\otimes\FrameBasis
+\FrameBasis\otimes\bm\tau+\bm\tau\otimes\FrameBasis
+\StressTensor_{\Section},
\]
where \(\StressTensor_{\Section}:=(\PlaneBasis_\alpha\otimes\PlaneBasis_\beta)\,\StressTensor\) is \emph{a priori} nonzero.  
Ie\c{s}an’s construction shows that, for the Saint–Venant class recovered below, the in–plane block vanishes, i.e.
\begin{equation}
\StressTensor=\AxialStress\,\FrameBasis\otimes\FrameBasis
+\FrameBasis\otimes\bm{\tau}+\bm{\tau}\otimes\FrameBasis,
\label{eq:saint-venant-stress}
\end{equation}
so the ``Saint–Venant-type stress hypothesis" is an \emph{outcome} of the solution.

\subsection{Problem 1: extension--bending--torsion}

\subsubsection*{Statement}
Given resultants on \(\Section_1\) with \(F_\alpha=0\) (\(\alpha=1,2\)) and prescribed \(N=F_3\), \(\bm M=(M_1,M_2,M_3)\),
find \((\bm u,\StressTensor)\) such that:
\begin{itemize}
\item \(\bm u,\StressTensor\) satisfy \eqref{eq:partial-elastostatic-equilibrium} in \(\Body\);
\item end tractions on \(\Section_0,\Section_L\) are unknown pointwise with resultants producing the given \(N,\bm M\) via \eqref{eq:def-resultants};
\item \textbf{Ie\c{s}an’s kinematic postulate:} the axial derivative is rigid on each section,
\(
\bm u'=\bm\alpha+\bm\beta\times\bm{r}
\),
with constants \(a_\varepsilon=\FrameBasis\cdot\bm\alpha\),
\(a_\vartheta=\FrameBasis\cdot\bm\beta\),
\(\bm a_\Section=\FrameBasis\times\bm\beta_\kappa\).
\end{itemize}
Under this postulate the torsion warping \(\WarpTwistIesan\) is defined by the Neumann problem
\begin{equation}
\Delta\WarpTwistIesan=0\ \text{ in }\Section,\qquad
\partial_{\bm\nu}\WarpTwistIesan=\varepsilon_{\alpha\beta}n_\alpha x_\beta\ \text{ on }\SectionBoundary.
\label{eq:bvp-warp-twist}
\end{equation}

\subsubsection*{Solution}
The sectional displacement is
\begin{equation}
\begin{aligned}
\hat{\bm u}\{\bm a\}
&=\Big(-\tfrac12 x^2\,\mathbf{1}-\nu\,\operatorname{dev}\bm{r}\otimes\bm{r}+x\,\FrameBasis\otimes\bm{r}\Big)\bm a_\Section +(x\,\FrameBasis-\nu\,\bm{r})\,a_\varepsilon
+\big(x\,\FrameBasis\times\bm{r}+\WarpTwistIesan\,\FrameBasis\big)a_\vartheta ,
\end{aligned}
\label{eq:iesan-displacement-1}
\end{equation}
with strain
\begin{equation}
\bm E^{(a)}=
\big(a_\varepsilon+\bm{r}\cdot\bm a_\Section\big)\,(\FrameBasis\otimes\FrameBasis-\nu\,\mathbf{1}_\Section)
+\tfrac12\big(\FrameBasis\otimes\bm\varepsilon_\vartheta+\bm\varepsilon_\vartheta\otimes\FrameBasis\big),
\quad
\bm\varepsilon_\vartheta:=a_\vartheta\big(\FrameBasis\times\bm{r}+\nabla\WarpTwistIesan\big),
\label{eq:iesan-strain-tensor-1}
\end{equation}
and axial/shear components
\begin{equation}
\bm{\tau}=G\big(\nabla\WarpTwistIesan-\FrameBasis\times\bm{r}\big)a_\vartheta,
\qquad
\AxialStress=E\big(\bm{r}\cdot\bm a_\Section+a_\varepsilon\big).
\label{eq:iesan-stress-1}
\end{equation}
The parameters \(\bm a=(\bm a_\Section,a_\varepsilon,a_\vartheta)\) are fixed by the resultants:
\begin{equation}
E\big(\IntRoR\,\bm a_\Section+A\,a_\varepsilon\,\CentroidLocation\big)=\FrameBasis\times\bm M,\qquad
EA\big(\bm a_\Section\cdot\CentroidLocation+a_\varepsilon\big)=-N,\qquad
G\,\Jo\,a_\vartheta=-T,
\label{eq:2-15}
\end{equation}
with
\(
\IntRoR=\int_\Section \bm{r}\otimes\bm{r}\,\mathrm{d}A
\)
and torsion constant
\begin{equation}
\Jo=\int_\Section\Big(r^2+(\FrameBasis\times\bm{r})\cdot\nabla\WarpTwistIesan\Big)\,\mathrm{d}A .
\label{eq:def-j-origin}
\end{equation}
Within this recovered solution one has \(\StressTensor_{\Section}=\mathbf{0}\), i.e. the Saint–Venant form \eqref{eq:saint-venant-stress}.

\subsection{Problem 2: flexure}
\subsubsection*{Statement}
Given a transverse shear resultant \(\ShearForce\) on \(\Section_1\) with
\(F_3=0\) and \(\bm M=\mathbf{0}\), find an elastostatic field \((\bm u,\StressTensor)\) such that:
\begin{itemize}
\item \(\bm u,\StressTensor\) satisfy \eqref{eq:partial-elastostatic-equilibrium};
\item end tractions on \(\Section_0,\Section_L\) are unknown pointwise but their resultants equal the given \(\ShearForce\) and zero moments via \eqref{eq:def-resultants};
\item \textbf{Ie\c{s}an’s kinematic postulate:} the axial derivative is rigid on each section,
\(
\bm u'=\bm\alpha+\bm\beta\times\bm{r}
\)
(as above), which ensures uniqueness in the relaxed class.
\end{itemize}

\subsubsection*{Solution}
The flexure parameters are determined by
\begin{equation}
\begin{aligned}
E\big(\IntRoR\,\bm b_\Section+A\,b_\varepsilon\,\CentroidLocation\big)&=-\ShearForce,\\
\bm b_\Section\cdot\CentroidLocation+b_\varepsilon&=0,\\
b_\vartheta&=0,
\end{aligned}
\label{eq:2-18}
\end{equation}
and the flexural warping \(\WarpShearIesan\) solves
\begin{equation}
\left\{\begin{aligned}
\Delta\WarpShearIesan&=-2\big((\bm{r}-\CentroidLocation)\cdot\bm b_\Section\big)\quad &&\text{in }\Section,\\
\partial_{\bm\nu}\WarpShearIesan&=\nu\,\bm\nu\cdot\Big(\operatorname{dev}\bm{r}\otimes\bm{r}-\bm{r}\otimes\CentroidLocation\Big)\bm b_\Section
\quad &&\text{on }\SectionBoundary,\\
\int_\Section \WarpShearIesan\,\mathrm{d}A&=0. &&
\end{aligned}\right.
\label{eq:bvp-warp-shear-iesan}
\end{equation}
Up to rigid motion, the displacement is
\begin{equation}
\begin{aligned}
\hat{\bm u}_\Section(x,\bm{r})
&=-\tfrac16\,x^3\,\bm b_\Section
+x\,\nu\Big(\tfrac12 r^2\,\bm b_\Section-(\bm b_\Section\cdot\bm{r})\,\bm{r}\Big)
+x\,\nu\,(\CentroidLocation\cdot\bm b_\Section)\,\bm{r}
+c_\vartheta\,x\,(\FrameBasis\times\bm{r}),\\
\FrameBasis\cdot\hat{\bm u}(x,\bm{r})
&=\tfrac12\,x^2\,\big((\bm{r}-\CentroidLocation)\cdot\bm b_\Section\big)+c_\vartheta\,\WarpTwistIesan(\bm{r})+\WarpShearIesan(\bm{r}),
\end{aligned}
\label{eq:iesan-displacement-2}
\end{equation}
with \(c_\vartheta\) fixed by
\begin{equation}
\Jo\,c_\vartheta
=-\int_\Section\Big[(\FrameBasis\times\bm{r})\cdot\nabla\big(\bm{\WarpShearIesan}\cdot\bm b_\Section\big)
+\tfrac12\,\nu\,r^2\,(\FrameBasis\times\bm{r})\cdot\bm b_\Section\Big]\mathrm{d}A .
\label{eq:iesan-2-22}
\end{equation}
The strain and the Saint–Venant axial/shear components are
\begin{equation}
\begin{aligned}
\bm E(\hat{\bm u})
&=(\varepsilon\,\FrameBasis+\bm\varepsilon_\gamma)\stackrel{\rm s}{\otimes}\FrameBasis
-\nu\,\varepsilon\,\mathbf{1}_\Section,\\
\varepsilon&=x\,(\bm{r}-\CentroidLocation)\cdot\bm b_\Section,\\
\bm\varepsilon_\gamma
&=(a_\vartheta+c_\vartheta)\big(\nabla\WarpTwistIesan+\FrameBasis\times\bm{r}\big)
-\nu\Big(\operatorname{dev}(\bm{r}\otimes\bm{r})-\bm{r}\otimes\CentroidLocation\Big)\bm b_\Section
+\nabla\big(\bm{\WarpShearIesan}\cdot\bm b_\Section\big),
\end{aligned}
\label{eq:iesan-strain-2}
\end{equation}
\begin{equation}
\AxialStress(\hat{\bm u})=E\,x\,(\bm{r}-\CentroidLocation)\cdot\bm b_\Section,\qquad
\bm{\tau}(\hat{\bm u})
=G\Big[(a_\vartheta+c_\vartheta)\big(\nabla\WarpTwistIesan+\FrameBasis\times\bm{r}\big)
-\nu\Big(\operatorname{dev}(\bm{r}\otimes\bm{r})-\bm{r}\otimes\CentroidLocation\Big)\bm b_\Section
+\nabla\big(\bm{\WarpShearIesan}\cdot\bm b_\Section\big)\Big].
\label{eq:iesan-stress-2}
\end{equation}
As in \(P_1\), the recovered field satisfies \(\StressTensor_{\Section}=\mathbf{0}\), hence \eqref{eq:saint-venant-stress}.

\subsection{General solution}
Combining \eqref{eq:iesan-displacement-1} and \eqref{eq:iesan-displacement-2} yields
\begin{equation}
\begin{aligned}
\hat{\bm u}\{\bm a,\bm b_\Section\}
&=\Big(-\tfrac12 x^2\,\mathbf{1}-\nu\,\operatorname{dev}\bm{r}\otimes\bm{r}+x\,\FrameBasis\otimes\bm{r}\Big)\bm a_\Section
+(x\,\FrameBasis-\nu\,\bm{r})\,a_\varepsilon
+\big(x\,\FrameBasis\times\bm{r}+\WarpTwistIesan\,\FrameBasis\big)a_\vartheta\\
&\quad-\tfrac16\,x^3\,\bm b_\Section
-x\,\nu\Big(\operatorname{dev}\bm{r}\otimes\bm{r}-\bm{r}\otimes\CentroidLocation\Big)\bm b_\Section
+\tfrac12 x^2\,\FrameBasis\otimes(\bm{r}-\CentroidLocation)\,\bm b_\Section
+\FrameBasis\otimes\bm{\WarpShearIesan}(\bm{r})\,\bm b_\Section\\
&\quad+\big(x\,\FrameBasis\times\bm{r}+\WarpTwistIesan\,\FrameBasis\big)c_\vartheta,
\end{aligned}
\label{eq:iesan-general-solution}
\end{equation}
with \(\bm a\) fixed by \eqref{eq:2-15}, \(\bm b\) by \eqref{eq:2-18}, \(\WarpTwistIesan\) by \eqref{eq:bvp-warp-twist},
\(\WarpShearIesan\) by \eqref{eq:bvp-warp-shear-iesan}, and \(c_\vartheta\) by \eqref{eq:iesan-2-22}.


\subsubsection{Equilibrium in the Saint-Venant class}
For any superposition satisfying \(\StressTensor\,\bm n=\mathbf{0}\) on \(\BodyLateralSurface\) and \eqref{eq:saint-venant-stress},
the field equations reduce on each section to
\begin{subequations}
\label{eq:sv-equilibrium}
\begin{align}
\bm{\tau}'&=\mathbf{0} && \text{in }\Section,\\
\DivS\bm{\tau}+\AxialStress'&=0 && \text{in }\Section,\\
\bm{\tau}\cdot\bm\nu&=0 && \text{on }\SectionBoundary.
\end{align}
\end{subequations}
In particular, \(\bm{\tau}=\bm{\tau}(\bm{r})\) is \(x\)-independent, and for flexure
\(\DivS\bm{\tau}=-E\,(\bm{r}-\CentroidLocation)\cdot\bm b_\Section\).
